% titel: Blatt 7, A4: Lewis-Strukturen und pi-Delokalisierung
% skript: Mortimer Kap. 9 und 10 (Lewis, delokalisierte pi-Systeme)

% Teil I: Werkzeuge, kein L\"osungsweg.

Zu zeichnen sind Lewis-Strukturformeln f\"ur $\mathrm{NO_3^-}$,
$\mathrm{H_2CO}$, $\mathrm{CO_3^{2-}}$, $\mathrm{O_3}$, $\mathrm{NH_4^+}$
und Benzol ($\mathrm{C_6H_6}$), jeweils mit der Frage nach
$\pi$-delokalisierten Bindungen.

\begin{werkzeugkasten}[pi-Delokalisierung erkennen]
\begin{itemize}
  \item \textbf{Delokalisiert:} mehr als eine gleichwertige
        Resonanzstruktur m\"oglich; das Elektron liegt in einem
        $\pi$-System \"uber mehr als zwei Atome verteilt
        (Mesomerie/Resonanz).
  \item Typische F\"alle: $\mathrm{NO_3^-}$ (3 mesomere Formen),
        $\mathrm{CO_3^{2-}}$ (3), $\mathrm{O_3}$ (2), Benzol (2
        Kekul\'e-Formen mit delokalisiertem $\pi$-Ring).
  \item \textbf{Lokalisiert:} nur eine Lewis-Struktur ohne
        Resonanz, z.~B. einfache C=O-Doppelbindung ($\mathrm{H_2CO}$) oder
        rein ionisches $\mathrm{NH_4^+}$.
\end{itemize}
\textit{Skript:} Mortimer Kap. 10, Resonanz/Delokalisierung.
\end{werkzeugkasten}

\begin{achtung}[Darauf achten]
\begin{itemize}
  \item Ladung z\"ahlen: $\mathrm{NO_3^-}$ und $\mathrm{CO_3^{2-}}$ sind
        Anionen (minuszeichen), $\mathrm{NH_4^+}$ ein Kation.
  \item Mehrere gleichwertige Resonanzstrukturen $\rightarrow$
        delokalisierte $\pi$-Bindungen; die Bindungsl\"angen werden dabei
        gleich lang.
  \item $\mathrm{H_2CO}$ und $\mathrm{NH_4^+}$ besitzen nur eine
        dominante Lewis-Struktur $\rightarrow$ keine nennenswerte
        $\pi$-Delokalisierung.
\end{itemize}
\end{achtung}
